\documentclass[10pt,a4paper]{amsart}
\usepackage{amsmath,amssymb,amsthm,graphicx}
% Retained minimal preamble for the Stacks Project chapter sites-cohomology.tex (Cohomology on Sites).
% Verbatim excerpt of the chapter preamble: Xy-pic setup (native preamble lines 14-18) and the native
% theorem declarations with the mathematical macros (native preamble lines 154-217). Copied mechanically;
% nothing added, removed or redefined. The wrapper class and the amsmath/amssymb/amsthm/graphicx base
% packages are supplied by the assembler, so the class/externaldocument part of the native preamble is omitted.
% For commutative diagrams we use Xy-pic
\usepackage[all]{xy}

% We use 2cell for 2-commutative diagrams.
\xyoption{2cell}
% Theorem environments.
%
\theoremstyle{plain}
\newtheorem{theorem}[subsection]{Theorem}
\newtheorem{proposition}[subsection]{Proposition}
\newtheorem{lemma}[subsection]{Lemma}

\theoremstyle{definition}
\newtheorem{definition}[subsection]{Definition}
\newtheorem{example}[subsection]{Example}
\newtheorem{exercise}[subsection]{Exercise}
\newtheorem{situation}[subsection]{Situation}

\theoremstyle{remark}
\newtheorem{remark}[subsection]{Remark}
\newtheorem{remarks}[subsection]{Remarks}

\numberwithin{equation}{subsection}

% Macros
%
\def\lim{\mathop{\mathrm{lim}}\nolimits}
\def\colim{\mathop{\mathrm{colim}}\nolimits}
\def\Spec{\mathop{\mathrm{Spec}}}
\def\Hom{\mathop{\mathrm{Hom}}\nolimits}
\def\Ext{\mathop{\mathrm{Ext}}\nolimits}
\def\SheafHom{\mathop{\mathcal{H}\!\mathit{om}}\nolimits}
\def\SheafExt{\mathop{\mathcal{E}\!\mathit{xt}}\nolimits}
\def\Sch{\mathit{Sch}}
\def\Mor{\mathop{\mathrm{Mor}}\nolimits}
\def\Ob{\mathop{\mathrm{Ob}}\nolimits}
\def\Sh{\mathop{\mathit{Sh}}\nolimits}
\def\NL{\mathop{N\!L}\nolimits}
\def\CH{\mathop{\mathrm{CH}}\nolimits}
\def\proetale{{pro\text{-}\acute{e}tale}}
\def\etale{{\acute{e}tale}}
\def\QCoh{\mathit{QCoh}}
\def\Ker{\mathop{\mathrm{Ker}}}
\def\Im{\mathop{\mathrm{Im}}}
\def\Coker{\mathop{\mathrm{Coker}}}
\def\Coim{\mathop{\mathrm{Coim}}}

% Boxtimes
%
\DeclareMathSymbol{\boxtimes}{\mathbin}{AMSa}{"02}

%
% Macros for moduli stacks/spaces
%
\def\QCohstack{\mathcal{QC}\!\mathit{oh}}
\def\Cohstack{\mathcal{C}\!\mathit{oh}}
\def\Spacesstack{\mathcal{S}\!\mathit{paces}}
\def\Quotfunctor{\mathrm{Quot}}
\def\Hilbfunctor{\mathrm{Hilb}}
\def\Curvesstack{\mathcal{C}\!\mathit{urves}}
\def\Polarizedstack{\mathcal{P}\!\mathit{olarized}}
\def\Complexesstack{\mathcal{C}\!\mathit{omplexes}}
% \Pic is the operator that assigns to X its picard group, usage \Pic(X)
% \Picardstack_{X/B} denotes the Picard stack of X over B
% \Picardfunctor_{X/B} denotes the Picard functor of X over B
\def\Pic{\mathop{\mathrm{Pic}}\nolimits}
\def\Picardstack{\mathcal{P}\!\mathit{ic}}
\def\Picardfunctor{\mathrm{Pic}}
\def\Deformationcategory{\mathcal{D}\!\mathit{ef}}
\usepackage{cleveref}
\crefname{originallemma}{original lemma}{original lemmas}
\Crefname{originallemma}{Original lemma}{Original lemmas}
\makeatletter
\newlabel{lemma-diagonal-better}{{0BYT}{source}{}{sourcepointer-1}{}}
\newlabel{lemma-diagonal-better@cref}{{[originallemma][0][]0BYT}{[1][0][]source}{}{}{}}
\makeatother
\title{Canonical morphism from a derived tensor product into an internal hom}
\author{The Stacks Project Authors}
\date{stacks-project:sites-cohomology.tex:local-snapshot-20261011; GNU FDL 1.2 or later}
\begin{document}
\maketitle
\noindent\textit{Editorial note.} All delivered mathematical text is the author's own, copied from the pinned Stacks Project chapter file sites-cohomology.tex; no statement or proof step was rewritten, reordered or removed.\par\medskip
\noindent\textit{Editorial note.} The proof is the author's complete construction of the asserted morphism: the representatives, the quasi-isomorphism into a K-injective complex and the composite of the two displayed arrows are all in the counted span.\par\medskip
\noindent\textit{Editorial note.} Cross-chapter restatement: the same label occurs in chapter cohomology and in chapter more-algebra. Those statements are the special cases of this ringed-site statement, so per the lane-s4 rule only this more general occurrence is counted here.\par\medskip
\noindent\textit{Editorial note.} The single labelled reference occurring inside the counted spans is carried in assembly.external\_references as a non-load-bearing navigation pointer to the same chapter.\par\medskip
\noindent\textit{Editorial note.} The preamble is the chapter's own retained minimal preamble (sources/stacks-sites-cohomology/retained-preamble.tex), a verbatim excerpt of the native preamble: the Xy-pic setup and the native theorem declarations with the mathematical macros.\par\medskip
\noindent\textit{Allowed standard prerequisites (editorial).}\begin{itemize}
\item K-flat and K-injective complexes of O-modules. Same chapter, the definitions of K-flat and K-injective complexes and the existence of quasi-isomorphisms into a K-injective complex. Supplies the representatives K\textasciicircum{}bullet, I\textasciicircum{}bullet, J\textasciicircum{}bullet chosen at the start of the counted proof.
\item RSheafHom on the derived category of a ringed site. Same chapter, the derived internal hom represented by the complex of sheaves of homomorphisms into a K-injective complex. Supplies the target of the constructed morphism and the identification of SheafHom\textasciicircum{}bullet(M\textasciicircum{}bullet, J\textasciicircum{}bullet) with RSheafHom(M, K tensor\textasciicircum{}L L).
\item Derived tensor product computed by a K-flat representative. Same chapter, the derived tensor product and the role of K-flatness in computing it. Supplies the source K tensor\textasciicircum{}L RSheafHom(M, L) as represented by the total complex in the display.
\end{itemize}\medskip
\ifdefined\hypertarget\hypertarget{sourcepointer-1}{}\fi
\noindent\textit{Source reference.} Original lemma 0BYT: sources/stacks-sites-cohomology/pinned-original-source.tex, source lines 8659--8659. Same-chapter navigation pointer invoked at proof line 9083, where the author takes the first of the two displayed maps from that earlier lemma. The pointer only resolves the reference for the reader; the cited lemma is not reproduced inside the counted span, and no bibliography citation is involved.\par\medskip
\begin{lemma}
\label{lemma-internal-hom-diagonal-better}
Let $(\mathcal{C}, \mathcal{O})$ be a ringed site. Given $K, L, M$
in $D(\mathcal{O})$ there is a canonical morphism
$$
K \otimes_\mathcal{O}^\mathbf{L} R\SheafHom(M, L)
\longrightarrow
R\SheafHom(M, K \otimes_\mathcal{O}^\mathbf{L} L)
$$
in $D(\mathcal{O})$ functorial in $K, L, M$.
\end{lemma}

\begin{proof}
Choose a K-flat complex $\mathcal{K}^\bullet$ representing $K$,
and a K-injective complex $\mathcal{I}^\bullet$ representing $L$, and
choose any complex of $\mathcal{O}$-modules $\mathcal{M}^\bullet$
representing $M$. Choose a quasi-isomorphism
$\text{Tot}(\mathcal{K}^\bullet \otimes_{\mathcal{O}_X} \mathcal{I}^\bullet)
\to \mathcal{J}^\bullet$
where $\mathcal{J}^\bullet$ is K-injective. Then we use the map
$$
\text{Tot}\left(
\mathcal{K}^\bullet \otimes_\mathcal{O}
\SheafHom^\bullet(\mathcal{M}^\bullet, \mathcal{I}^\bullet)
\right)
\to
\SheafHom^\bullet(\mathcal{M}^\bullet,
\text{Tot}(\mathcal{K}^\bullet \otimes_\mathcal{O} \mathcal{I}^\bullet))
\to
\SheafHom^\bullet(\mathcal{M}^\bullet, \mathcal{J}^\bullet)
$$
where the first map is the map from Lemma \ref{lemma-diagonal-better}.
\end{proof}

\end{document}
